\documentclass[11pt]{article}

\usepackage[a4paper,margin=28mm]{geometry}
\usepackage{amsmath,amssymb,amsthm,mathtools}
\usepackage{graphicx}
\usepackage{enumitem}
\usepackage{microtype}
\usepackage[hidelinks]{hyperref}

\newtheorem{theorem}{Theorem}
\newtheorem{proposition}{Proposition}
\newtheorem{corollary}{Corollary}
\newtheorem{lemma}{Lemma}
\newtheorem{definition}{Definition}
\newtheorem{remark}{Remark}

\newcommand{\N}{\mathbb N}
\newcommand{\eps}{\varepsilon}
\newcommand{\conv}{\operatorname{conv}}
\newcommand{\Def}{\operatorname{Def}}
\newcommand{\norm}[1]{\lVert #1\rVert}

\title{A Strengthened Tree Banach--Saks Partial Answer to Braga's WSC Coarse Embedding Problem}
\author{FA/Banach discovery run \texttt{fa\_banach\_001}}
\date{Updated June 22, 2026}

\begin{document}
\maketitle

\section*{Status}

This packet concerns Bruno de Mendonca Braga's paper
\emph{On asymptotically uniformly smoothness and nonlinear geometry of Banach
spaces}, arXiv:1808.03254 \cite{BragaJIMJ}.  It incorporates the supplied
report
\[
\texttt{evidence/supplied\_braga\_weak\_p\_banach\_saks\_note\_2026\_06\_22/}.
\]
The result is a strengthened partial answer, likely valid pending human
review.  It is not a full solution of Problem 8.6.

\section*{Source Problem}

Problem 8.6 of the source paper asks:

\begin{quote}
Let \(X\) and \(Y\) be Banach spaces. Assume that \(Y\) has the weak
\(p\)-Banach--Saks property, for some \(p\in(1,\infty)\). If \(X\) coarse
Lipschitzly embeds into \(Y\) by a weakly sequentially continuous map, does it
follow that \(X\) has the weak \(p\)-Banach--Saks property?
\end{quote}

\begin{center}
\includegraphics[width=\textwidth]{figures/open_problem_crop.png}
\end{center}

\section*{Classification}

The packet proves positive results under stronger or additional hypotheses.
The supplied report also explains why the original weak-\(p\) hypothesis does
not feed directly into the nonlinear telescoping proof: the construction
creates weakly null trees, while weak \(p\)-Banach--Saks controls only one
weakly null sequence at a time.  Braga's own examples show that this
tree-versus-sequence gap is genuine \cite[Theorem 7.1]{BragaJIMJ}.

The previous packet's tree-\(p\) permanence theorem is correct.  The report
adds three local repairs: the signed stabilization citation is Braga's
Proposition 6.3 rather than Proposition 6.5; normalization of semi-normalized
weakly null increments preserves weak nullity; and the rescaled translated map
used in the proof remains weakly sequentially continuous.

\section*{Definitions}

For a map \(f:X\to Y\) between Banach spaces, set
\[
 \omega_f(t)=\sup\{\norm{f(x)-f(y)}:\norm{x-y}\le t\},
 \quad
 \rho_f(t)=\inf\{\norm{f(x)-f(y)}:\norm{x-y}\ge t\}.
\]
The map is a coarse embedding if \(\omega_f(t)<\infty\) for all \(t\) and
\(\rho_f(t)\to\infty\).  It is coarse Lipschitz if there are \(a,b>0\) and
\(A\ge0\) such that
\begin{equation}\label{eq:affine-control}
 a\norm{x-y}-A\le\norm{f(x)-f(y)}\le b\norm{x-y}+A
 \qquad(x,y\in X).
\end{equation}
It is weakly sequentially continuous if \(x_n\rightharpoonup x\) implies
\(f(x_n)\rightharpoonup f(x)\).

A Banach space \(Z\) has weak \(p\)-Banach--Saks constant \(C\) if every
weakly null sequence \((z_n)\subset B_Z\) has an infinite \(M\subset\N\) such
that
\[
 \norm{\sum_{j=1}^k z_{n_j}}\le Ck^{1/p}
 \qquad(n_1<\cdots<n_k,\ n_j\in M).
\]
It has tree-\(p\)-Banach--Saks constant \(C\) if every normalized weakly null
tree of height \(k\) has a branch satisfying the same \(Ck^{1/p}\) upper
bound for the branch sum.  Tree-\(p\) implies weak \(p\) by applying the tree
property to the tree generated by a sequence \cite[Remark 6.4]{BragaJIMJ}.

\begin{proposition}[Signed stabilization]\label{prop:signed-tree}
If \(Z\) has tree-\(p\)-Banach--Saks constant \(C\), then every normalized
weakly null tree of height \(k\) has a branch on which all signed branch sums
have norm at most \(Ck^{1/p}\).
\end{proposition}

This is Braga's Proposition 6.3 \cite{BragaJIMJ}.  It follows by applying the
tree property to finitely many levelwise sign changes and then using Ramsey
stabilization.

\begin{lemma}[Sign-cube convexity]\label{lem:sign-cube}
If \(e_1,\ldots,e_k\) are vectors in a normed space and
\(t_1,\ldots,t_k\in[0,1]\), then
\[
 \sum_{j=1}^k t_j e_j\in
 \conv\{\sum_{j=1}^k\eps_j e_j:\eps_j\in\{-1,1\}\}.
\]
Consequently, if all signed sums on the right have norm at most \(K\), then
\(\norm{\sum_j t_j e_j}\le K\).
\end{lemma}

\begin{proof}
Choose independent signs \(\eps_j\) with expectation \(t_j\).  Then
\(\sum_j t_j e_j\) is the expectation of \(\sum_j\eps_j e_j\), hence lies in
the convex hull.
\end{proof}

\section*{Proof Intuition}

The nonlinear map is applied to partial sums along a weakly null tree.  The
increments
\[
 f(r(u_s+x_{s^\frown n}))-f(ru_s)
\]
are weakly null successor sequences because \(f\) is weakly sequentially
continuous.  Compression gives a lower bound on their norms and expansion
gives an upper bound.  After normalization, tree-\(p\)-Banach--Saks gives a
branch where all signed sums are small.  The sign-cube lemma then handles the
unequal increment lengths, and telescoping converts the increment sum into the
single endpoint difference controlled from below by \(\rho_f\).

\section*{Modulus Transfer}

\begin{theorem}[Modulus transfer]\label{thm:modulus-transfer}
Let \(1<p\le\infty\).  Suppose that \(Y\) has tree-\(p\)-Banach--Saks
constant \(C\), and let \(f:X\to Y\) be a weakly sequentially continuous
coarse embedding.  Fix \(r>0\) such that \(\rho_f(r)>0\).  Then every
normalized weakly null tree in \(X\) of height \(k\) has a branch
\((x_1,\ldots,x_k)\) satisfying
\begin{equation}\label{eq:master-transfer}
 \rho_f\!\left(r\norm{\sum_{j=1}^k x_j}\right)
 \le C\,\omega_f(r)k^{1/p},
\end{equation}
with the convention \(k^{1/\infty}=1\).
\end{theorem}

\begin{proof}
After translation assume \(f(0)=0\).  For a nonempty node \(s\), let
\[
 u_s=\sum_{t\preceq s}x_t,\qquad
 y_s=f(ru_s)-f(ru_{s^-}).
\]
For each fixed predecessor, the successor sequence of increments is weakly
null.  Since adjacent partial sums are distance \(r\) apart,
\[
 0<\rho_f(r)\le\norm{y_s}\le\omega_f(r)<\infty.
\]
Thus \(e_s=y_s/\norm{y_s}\) is a normalized weakly null tree.  Apply
Proposition~\ref{prop:signed-tree} to \((e_s)\).  On the selected branch, put
\[
 t_j=\frac{\norm{y_j}}{\omega_f(r)}\in[0,1].
\]
By Lemma~\ref{lem:sign-cube},
\[
 \norm{\sum_{j=1}^k y_j}
 =\omega_f(r)\norm{\sum_{j=1}^k t_j e_j}
 \le C\,\omega_f(r)k^{1/p}.
\]
The left side is
\[
 \norm{f(r\sum_{j=1}^k x_j)-f(0)}.
\]
The definition of \(\rho_f\) gives \eqref{eq:master-transfer}.
\end{proof}

\begin{corollary}[Coarse Lipschitz tree permanence]\label{cor:CL-tree-transfer}
Let \(1<p<\infty\).  If \(Y\) has tree-\(p\)-Banach--Saks constant \(C\) and
\(f:X\to Y\) is weakly sequentially continuous and satisfies
\eqref{eq:affine-control}, then \(X\) has tree-\(p\)-Banach--Saks.  More
precisely, for every \(r>A/a\), every normalized weakly null tree has a branch
with
\[
 \norm{\sum_{j=1}^k x_j}
 \le \frac{C(br+A)+A}{ar}k^{1/p}.
\]
Letting \(r\to\infty\), the constants can be taken arbitrarily close to
\(bC/a\).  Hence \(X\) also has weak \(p\)-Banach--Saks.
\end{corollary}

\begin{proof}
Use \(\omega_f(r)\le br+A\), and apply the lower affine estimate in
\eqref{eq:affine-control} to the endpoint \(r\sum_jx_j\).  Since
\(k^{1/p}\ge1\), the additive \(A\) is absorbed as displayed.  The final
claim follows from tree-\(p\Rightarrow\) weak-\(p\).
\end{proof}

\begin{corollary}[Power compression]\label{cor:power-compression}
Under the hypotheses of Theorem~\ref{thm:modulus-transfer}, suppose
\(\rho_f(t)\ge ct^\alpha-d\) for some \(c>0\), \(d\ge0\), and \(\alpha>0\).
Then every normalized weakly null tree in \(X\) has a branch satisfying
\[
 \norm{\sum_{j=1}^k x_j}
 \le
 \frac{(C\omega_f(r)+d)^{1/\alpha}}{r c^{1/\alpha}}
 k^{1/(\alpha p)}.
\]
If \(\alpha p>1\), this gives tree-\((\alpha p)\)-Banach--Saks in \(X\).
\end{corollary}

\begin{remark}[Endpoint]
For \(p=\infty\), Theorem~\ref{thm:modulus-transfer} gives a branch whose
sums are uniformly bounded independently of the height.  Thus
tree-\(\infty\)-Banach--Saks passes under weakly sequentially continuous
coarse embeddings.
\end{remark}

\section*{Weak p Under Additive Defect}

The original weak-\(p\) hypothesis can be used directly if the embedding is
sufficiently additive on finite sums.  After translating \(f\), write
\(\widetilde f=f-f(0)\), and define
\[
 \Def_f(k,r)=\sup_{x_1,\ldots,x_k\in B_X}
 \norm{
  \widetilde f(r\sum_{j=1}^k x_j)
  -\sum_{j=1}^k\widetilde f(rx_j)
 }.
\]

\begin{theorem}[Defect criterion]\label{thm:defect}
Let \(1<p<\infty\).  Suppose that \(Y\) has weak \(p\)-Banach--Saks constant
\(C\), and let \(f:X\to Y\) be weakly sequentially continuous and satisfy
\eqref{eq:affine-control}.  If there are \(D\ge0\) and scales \(r_k\ge1\)
such that
\[
 \frac{\Def_f(k,r_k)}{r_k}\le Dk^{1/p}\qquad(k\in\N),
\]
then \(X\) has weak \(p\)-Banach--Saks, with constant at most
\[
 \frac{C(b+A)+D+A}{a}.
\]
\end{theorem}

\begin{proof}
Let \((x_n)\subset B_X\) be weakly null and fix \(k\).  Define
\[
 g_k(x)=\frac{\widetilde f(r_kx)}{r_k}.
\]
Then \(g_k\) is weakly sequentially continuous, \(g_k(0)=0\), and
\(\norm{g_k(x_n)}\le b+A\).  Apply weak \(p\)-Banach--Saks in \(Y\) to
\((g_k(x_n))\), obtaining an infinite \(M\) such that
\[
 \norm{\sum_{j=1}^k g_k(x_{n_j})}
 \le C(b+A)k^{1/p}
 \qquad(n_1<\cdots<n_k,\ n_j\in M).
\]
The defect assumption gives
\[
 \norm{g_k(\sum_{j=1}^k x_{n_j})}
 \le (C(b+A)+D)k^{1/p}.
\]
The lower affine estimate for \(g_k\), with \(A/r_k\le Ak^{1/p}\), gives the
displayed constant.
\end{proof}

\begin{corollary}[Bounded or sublinear additive defect]
The conclusion of Theorem~\ref{thm:defect} holds if either
\[
 \sup_{u,v\in X}\norm{f(u+v)-f(u)-f(v)+f(0)}<\infty
\]
or, with
\[
 \delta_f(R)=\sup_{\norm u,\norm v\le R}
 \norm{f(u+v)-f(u)-f(v)+f(0)},
\]
one has \(\delta_f(R)/R\to0\).
\end{corollary}

\begin{proof}
Bounded binary defect gives \(\Def_f(k,r)\le(k-1)D_0\); take \(r_k=k\).  For
sublinear defect, use
\[
 \Def_f(k,r)\le(k-1)\delta_f(kr)
\]
and choose \(r_k\) so large that the ratio is at most \(k^{1/p}\).
\end{proof}

\section*{Why This Still Does Not Settle Problem 8.6}

Weak sequential continuity makes successors of each fixed tree node weakly
null, but it does not make an enumeration of all increments weakly null,
because the base point changes from node to node.  Weak \(p\)-Banach--Saks can
control a single sequence such as \(f(rx_n)-f(0)\); the coarse lower estimate
instead sees \(f(r\sum_jx_{n_j})-f(0)\).  The difference between these two
quantities is exactly the additive defect isolated above.

The supplied report also checked natural counterexample routes.  The reflexive
Kalton--Peck space \(Z_p\) remains a possible source obstruction, but no
suitable weakly sequentially continuous embedding into a weak-\(p\) target was
found.  Braga's Baire-sum construction gives weak-\(p\) spaces with very bad
trees, but no coarse Lipschitz weakly sequentially continuous embedding
realizing the bad tree behavior was obtained.  These failures are not
counterexamples.

\section*{Literature and Novelty Status}

The supplied report records searches through June 21, 2026 for the exact
problem number, weak-\(p\) Banach--Saks phrases, weakly sequentially continuous
coarse Lipschitz embeddings, and forward citations of the source.  Later
nearby work develops tree-\(p\)-Banach--Saks and Hamming concentration
properties \cite{FovelleHamming}, studies asymptotic smoothness and strict
hierarchies \cite{CauseyFovelleLancien}, and gives further concentration
examples \cite{FovelleSmoothness}.  No located paper states a proof or a
counterexample to Problem 8.6.

Thus the packet should be reviewed as a strengthened partial candidate, not as
a full solution.

\begin{thebibliography}{99}

\bibitem{BaudierLancienBook}
F. P. Baudier and G. Lancien,
\emph{Asymptotic and nonlinear geometries of Banach spaces and their interactions},
arXiv:2512.00817, 2025.

\bibitem{BragaWSC}
B. de Mendonca Braga,
\emph{Nonlinear weakly sequentially continuous embeddings between Banach spaces},
Int. Math. Res. Not. IMRN (2020), no. 18, 5506--5533.

\bibitem{BragaJIMJ}
B. de Mendonca Braga,
\emph{On asymptotic uniform smoothness and nonlinear geometry of Banach spaces},
J. Inst. Math. Jussieu \textbf{20} (2021), no. 1, 65--102;
arXiv:1808.03254.

\bibitem{CauseyFovelleLancien}
R. M. Causey, A. Fovelle, and G. Lancien,
\emph{Asymptotic smoothness in Banach spaces, three-space properties and coarse Lipschitz geometry},
arXiv:2110.06710, version 2, 2022.

\bibitem{FovelleHamming}
A. Fovelle,
\emph{Hamming graphs and concentration properties in Banach spaces},
arXiv:2106.04297, version 3, 2023.

\bibitem{FovelleSmoothness}
A. Fovelle,
\emph{Asymptotic smoothness, concentration properties in Banach spaces and applications},
arXiv:2311.15900, version 2, 2024.

\end{thebibliography}

\end{document}
